% =============================================================================
% Elementos pré-textuais (ABNT NBR 14724):
% capa, folha de rosto, resumo, abstract, lista de abreviaturas.
% =============================================================================

% ------------------------------- CAPA ---------------------------------------
\begin{titlingpage}
\begin{center}
  \textbf{\MakeUppercase{Instituto Federal Fluminense}}\\[2\baselineskip]
  \textbf{\MakeUppercase{William da Silva Vianna}}\\[6\baselineskip]
  \textbf{mcp-kicad-vscode: um servidor Model Context Protocol para automação
  do fluxo de projeto de placas de circuito impresso no KiCad}\\[10\baselineskip]
  \vfill
  Campos dos Goytacazes, RJ\\[0.5\baselineskip]
  2026
\end{center}
\end{titlingpage}

% --------------------------- FOLHA DE ROSTO ---------------------------------
\begin{titlingpage}
\begin{center}
  \textbf{\MakeUppercase{William da Silva Vianna}}\\[6\baselineskip]
  \textbf{mcp-kicad-vscode: um servidor Model Context Protocol para automação
  do fluxo de projeto de placas de circuito impresso no KiCad}\\[4\baselineskip]
\end{center}

\vspace{2\baselineskip}
\hfill
\begin{minipage}{0.55\textwidth}
  \SingleSpacing
  \footnotesize
  Trabalho apresentado ao Instituto Federal Fluminense (IFF) como requisito
  parcial para a obtenção do título de
  \textbf{[CURSO/TÍTULO A CONFIRMAR --- NÃO CONSTA NO WORKSPACE]}.

  Orientação: \textbf{[ORIENTADOR A CONFIRMAR --- NÃO CONSTA NO WORKSPACE]}.
\normalsize
\end{minipage}

\vfill
\begin{center}
  Campos dos Goytacazes, RJ\\
  2026
\end{center}
\end{titlingpage}

% -------------------------------- RESUMO ------------------------------------
\chapter*{Resumo}
\addcontentsline{toc}{chapter}{Resumo}

\noindent
O \emph{Model Context Protocol} (MCP) surgiu como um padrão aberto para
conectar assistentes baseados em grandes modelos de linguagem (LLMs) a
ferramentas e fontes de dados externas. Este trabalho especifica, implementa e
avalia o \texttt{mcp-kicad-vscode}, um servidor MCP dedicado à automação do
fluxo de projeto de placas de circuito impresso (PCB) no KiCad. A solução
adota arquitetura em duas camadas --- TypeScript para o protocolo e Python
para a interface com o KiCad --- e expõe 233 ferramentas em 24 módulos
funcionais, 173 delas indexadas para descoberta por palavra-chave em 16
categorias, além de 23 recursos dinâmicos de estado do projeto. O sistema
cobre esquemático, layout, roteamento, verificação de regras de projeto (DRC e
ERC), exportação de fabricação e integrações com os catálogos JLCPCB, LCSC e
DigiKey, com suporte ao roteador automático Freerouting. A avaliação combinou
métricas de implementação, execução das suítes de teste --- 84 casos
TypeScript e 2.367 casos Python coletados, com 2.328 aprovados, 8 falhas e 31
ignorados na execução local de setembro de 2026 --- e um estudo de caso com
uma placa seguidora de linhas de quatro camadas (97 \emph{footprints}, 83
redes, aproximadamente 82,1\,mm $\times$ 80,2\,mm), cuja verificação de regras
reportou 11 violações residuais (9 erros e 2 avisos) e nenhum item não
conectado. Discutem-se as limitações --- integração IPC não exercitada no
ambiente de validação, cobertura parcial de descoberta e avaliação restrita a
um único projeto --- e propõem-se trabalhos futuros.

\vspace{\baselineskip}
\noindent\textbf{Palavras-chave:} Model Context Protocol; KiCad; projeto de
placas de circuito impresso; automação de EDA; assistentes de inteligência
artificial.

% ------------------------------- ABSTRACT -----------------------------------
\chapter*{Abstract}
\addcontentsline{toc}{chapter}{Abstract}

\noindent
The \emph{Model Context Protocol} (MCP) has emerged as an open standard for
connecting assistants based on large language models (LLMs) to external tools
and data sources. This work specifies, implements and evaluates
\texttt{mcp-kicad-vscode}, an MCP server dedicated to automating the printed
circuit board (PCB) design workflow in KiCad. The solution adopts a two-layer
architecture --- TypeScript for the protocol and Python for the KiCad
interface --- and exposes 233 tools across 24 functional modules, 173 of which
are indexed for keyword discovery in 16 categories, plus 23 dynamic project
state resources. The system covers schematic capture, layout, routing, design
rule checking (DRC and ERC), fabrication export and integrations with the
JLCPCB, LCSC and DigiKey catalogs, with support for the Freerouting
autorouter. The evaluation combined implementation metrics, test suite
execution --- 84 TypeScript cases and 2,367 collected Python cases, with 2,328
passed, 8 failed and 31 skipped in the local run of September 2026 --- and a
case study with a four-layer line-follower board (97 footprints, 83 nets,
approximately 82.1\,mm $\times$ 80.2\,mm), whose design rule check reported 11
residual violations (9 errors and 2 warnings) and no unconnected items.
Limitations are discussed --- the IPC integration was not exercised in the
validation environment, discovery coverage is partial, and the evaluation is
restricted to a single project --- and future work is proposed.

\vspace{\baselineskip}
\noindent\textbf{Keywords:} Model Context Protocol; KiCad; printed circuit
board design; EDA automation; artificial intelligence assistants.

% -------------------- LISTA DE ABREVIATURAS E SIGLAS ------------------------
\chapter*{Lista de Abreviaturas e Siglas}
\addcontentsline{toc}{chapter}{Lista de Abreviaturas e Siglas}

\begin{tabular}{@{}p{2.6cm}p{11.4cm}@{}}
\toprule
\textbf{Sigla} & \textbf{Significado} \\
\midrule
ABNT & Associação Brasileira de Normas Técnicas \\
API & \emph{Application Programming Interface} \\
CI & \emph{Continuous Integration} \\
CLI & \emph{Command-Line Interface} \\
DRC & \emph{Design Rule Check} \\
EDA & \emph{Electronic Design Automation} \\
ERC & \emph{Electrical Rule Check} \\
IPC & \emph{Inter-Process Communication} (API IPC do KiCad) \\
JSON & \emph{JavaScript Object Notation} \\
JSON-RPC & \emph{JSON Remote Procedure Call} \\
LLM & \emph{Large Language Model} \\
LOC & \emph{Lines of Code} \\
MCP & \emph{Model Context Protocol} \\
MPPT & \emph{Maximum Power Point Tracking} \\
NDJSON & \emph{Newline-Delimited JSON} \\
PCB & \emph{Printed Circuit Board} \\
PTH & \emph{Plated Through-Hole} \\
SMD & \emph{Surface-Mount Device} \\
SWIG & \emph{Simplified Wrapper and Interface Generator} \\
USB & \emph{Universal Serial Bus} \\
\bottomrule
\end{tabular}
